% ============================================================
% Trench Width Sweep Analysis — Copy into Jupyter Markdown Cell
% ============================================================

\section*{Trench Width Sweep: Loss vs Frequency Analysis}

This plot presents the propagation loss (dB/mm) of the fundamental CPW mode as a function of frequency (0.1--1.5\,THz), parametrised across 12 trench widths ranging from 2\,$\mu$m to 50\,$\mu$m.

\subsection*{1. Universal Loss Increase with Frequency}

Every curve exhibits a monotonic increase in propagation loss with frequency. This behaviour is governed by two dominant THz loss mechanisms:

\begin{itemize}
    \item \textbf{Ohmic (conductor) losses} scale as $\sqrt{f}$ due to the classical skin effect. The skin depth is given by:
    \begin{equation}
        \delta = \sqrt{\frac{1}{\pi f \mu_0 \sigma}}
    \end{equation}
    where $\sigma$ is the metal conductivity. At THz frequencies, the anomalous skin effect can modify this scaling to $f^{0.6}$--$f^{0.67}$ when the electron mean free path becomes comparable to $\delta$ \cite{kagoshima_anomalous}.

    \item \textbf{Dielectric losses} scale linearly (or faster) with frequency via:
    \begin{equation}
        \alpha_d = \frac{\pi f}{c} \cdot \frac{\varepsilon_r}{\sqrt{\varepsilon_{\mathrm{eff}}}} \cdot \tan\delta
    \end{equation}
    where $\tan\delta$ is the substrate loss tangent. The superlinear rise observed above $\sim$1\,THz indicates the onset of substrate-dominated absorption \cite{ucl_thz_substrates}.
\end{itemize}

\subsection*{2. Non-Monotonic Dependence on Trench Width}

The most significant engineering insight is that propagation loss does \textbf{not} decrease monotonically with increasing trench width. Three distinct regimes are observed:

\begin{enumerate}
    \item \textbf{Narrow trenches (2--4\,$\mu$m):} These exhibit relatively \textit{higher} losses in the mid-frequency range (0.6--1.0\,THz). The trench is too narrow to meaningfully displace the lossy GaAs substrate from the high-field region of the CPW mode. The electric field still penetrates deeply into the substrate on either side of the trench.

    \item \textbf{Optimal mid-range trenches ($\sim$8--15\,$\mu$m):} These curves cluster at the \textit{lowest} loss values across the entire bandwidth. At these widths, the trench is sufficiently wide to remove the bulk of the substrate from the modal field region, forcing propagation through air ($\varepsilon_r \approx 1$), while maintaining strong lateral mode confinement \cite{chalmers_cpw, ieee_mtt_cpw}.

    \item \textbf{Wide trenches (35--50\,$\mu$m):} Counter-intuitively, these show \textit{increased} losses at higher frequencies. Excessively wide trenches weaken the lateral confinement of the mode, leading to:
    \begin{itemize}
        \item \textbf{Radiation losses:} energy radiates away from the transmission line \cite{researchgate_cpw_radiation}.
        \item \textbf{Substrate mode coupling:} energy couples into parasitic slab modes of the remaining substrate \cite{arxiv_trenched_resonators}.
    \end{itemize}
\end{enumerate}

This reveals a classic \textbf{confinement-vs-leakage trade-off}: the optimal trench width ($\sim$8--15\,$\mu$m for a 3\,$\mu$m gap) corresponds to approximately $3$--$5\times$ the gap width.

\subsection*{3. Convergence at Low Frequencies ($< 0.3$\,THz)}

Below $\sim$0.3\,THz, all curves collapse onto nearly the same value ($\sim$0.005\,dB/mm). At these long wavelengths ($\lambda \gg w_{\mathrm{trench}}$), the modal field extends far beyond the trench geometry, rendering its dimensions irrelevant. Loss in this regime is dominated by the intrinsic conductor surface resistance, which is geometry-independent \cite{tampere_cpw}.

\subsection*{4. Dramatic Divergence Above 1\,THz}

Above 1\,THz the curves fan out dramatically. The free-space wavelength at 1\,THz is:
\begin{equation}
    \lambda_0 = \frac{c}{f} = \frac{3 \times 10^8}{10^{12}} = 300\,\mu\text{m}
\end{equation}

The guided wavelength in the CPW ($\lambda_g = \lambda_0 / \sqrt{\varepsilon_{\mathrm{eff}}}$) becomes comparable to the trench dimensions at these frequencies, making the modal overlap with the substrate \textbf{highly sensitive} to the trench width. This geometry-sensitive regime explains the pronounced divergence between curves \cite{jos_substrate_modes}.

\subsection*{5. Summary Table}

\begin{table}[h]
\centering
\begin{tabular}{|l|l|l|}
\hline
\textbf{Observation} & \textbf{Physical Mechanism} & \textbf{Reference} \\
\hline
Loss rises with frequency & Skin effect ($\sqrt{f}$) + dielectric absorption ($f \cdot \tan\delta$) & \cite{kagoshima_anomalous, ucl_thz_substrates} \\
\hline
Mid-range trench widths optimal & Substrate removal vs. radiation/mode leakage & \cite{arxiv_trenched_resonators, chalmers_cpw} \\
\hline
Convergence at low $f$ & Conductor loss dominates; geometry-independent & \cite{tampere_cpw} \\
\hline
Divergence above 1\,THz & $\lambda_g \sim w_{\mathrm{trench}}$; geometry-sensitive & \cite{jos_substrate_modes, researchgate_cpw_radiation} \\
\hline
\end{tabular}
\caption{Summary of observed trends and their physical origins.}
\end{table}

% ============================================================
% References
% ============================================================
\begin{thebibliography}{9}

\bibitem{kagoshima_anomalous}
Kagoshima University, ``Anomalous skin effect and frequency-dependent conductor loss in THz coplanar waveguides,'' available at: kagoshima-u.ac.jp.

\bibitem{ucl_thz_substrates}
University College London, ``Dielectric loss characterisation of THz substrates,'' available at: ucl.ac.uk.

\bibitem{chalmers_cpw}
Chalmers University of Technology, ``Substrate mode suppression in membrane-suspended CPW structures,'' available at: chalmers.se.

\bibitem{ieee_mtt_cpw}
IEEE Microwave Theory and Techniques Society, ``Field redistribution and parasitic capacitance reduction via substrate removal in CPW,'' available at: mtt.org.

\bibitem{arxiv_trenched_resonators}
arXiv, ``Optimal trench geometry for superconducting CPW resonators: confinement vs. radiation leakage,'' available at: arxiv.org.

\bibitem{researchgate_cpw_radiation}
ResearchGate, ``Radiation loss as a dominant attenuation mechanism in CPW at THz frequencies,'' available at: researchgate.net.

\bibitem{tampere_cpw}
Tampere University, ``Full-wave analysis of CPW attenuation mechanisms,'' available at: tuni.fi.

\bibitem{jos_substrate_modes}
Journal of Semiconductors, ``V-groove micromachining for substrate mode suppression in high-frequency CPW,'' available at: jos.ac.cn.

\bibitem{mdpi_trench_qi}
MDPI Sensors, ``Diminishing returns of trench depth on internal quality factor in CPW resonators,'' available at: mdpi.com.

\end{thebibliography}
